\documentclass[11pt]{article}
\usepackage[utf8]{inputenc}
\usepackage[spanish,es-nodecimaldot]{babel}
\usepackage[a4paper,margin=2.4cm]{geometry}
\usepackage{amsmath,booktabs,graphicx,siunitx,url}

\title{\textbf{Test 1: Sensibilidad a Malla, Equilibrio y Condiciones de Borde}}
\author{Alanj Rivera}
\date{\today}

\begin{document}
\maketitle

\section{Problema y Objetivo}
El objetivo es evaluar la credibilidad del método de elementos finitos (MEF) lineal 1D aplicado a una barra sometida a carga axial uniforme y puntual. Se analiza la diferencia entre el acuerdo nodal y la exactitud del campo continuo, encontrando el error máximo en puntos medios para mallas de 2, 4 y 8 . Además, se estudia el efecto estático y cinemático al introducir una restricción cinemática adicional $u(L)=0$.

\section{Datos, Supuestos y Método}
Se considera una barra elástica lineal con $L = \SI{2000}{\mm}$, $E = \SI{200000}{\N\per\mm\squared}$, $A = \SI{1000}{\mm\squared}$, carga puntual $P = \SI{20000}{\N}$ en $+x$ y carga distribuida $q = \SI{10}{\N\per\mm}$ en $+x$[cite: 1, 22]. La ecuación gobernante del problema de borde es:
\begin{equation}
-(EA u'(x))' = q(x), \quad \text{con } u(0) = 0 \quad \text{y} \quad EA u'(L) = P
\end{equation}
La solución analítica de referencia independiente obtenida por integración directa es[cite: 1]:
\begin{equation}
u_{\text{ref}}(x) = \frac{(P + qL)x - \frac{1}{2}q x^2}{EA}
\end{equation}
Para la aproximación MEF lineal, la solución entre nodos interpola linealmente, generando un error en puntos medios $x_m = (x_e + x_{e+1})/2$ que se cuantifica como $e_{\text{máx}} = \max |u_h(x_m) - u_{\text{ref}}(x_m)|$[cite: 1, 26].

\section{Resultados y Verificación}
En la Tabla~\ref{tab:resultados} se resume la respuesta del modelo para las tres mallas uniformes y el caso con condición de contorno alterada[cite: 1, 30].

\begin{table}[h!]
\centering
\caption{Resultados del análisis MEF para distintas mallas y condición alterada.}
\label{tab:resultados}
\small
\begin{tabular}{lcccccc}
\toprule
\textbf{Caso} & $n$ & $u(L)$ [\text{mm}] & $R_{\text{izq}}$ [\text{N}] & $R_{\text{der}}$ [\text{N}] & Error medio [\text{mm}] & Balance [\text{N}] \\
\midrule
M2 & 2 & 0.3000 & -40000.0 & 0.0 & 0.006250 & $0.0$ \\
M4 & 4 & 0.3000 & -40000.0 & 0.0 & 0.001563 & $2.18 \times 10^{-11}$ \\
M8 & 8 & 0.3000 & -40000.0 & 0.0 & 0.000391 & $-3.64 \times 10^{-11}$ \\
BC\_ALTERADO & 8 & 0.0000 & -10000.0 & -30000.0 & --- & $0.0$ \\
\bottomrule
\end{tabular}
\end{table}

\begin{figure}[h!]
\centering
\includegraphics[width=0.55\linewidth]{salidas/convergencia.png}
\caption{Convergencia del error absoluto máximo en puntos medios en escala log-log.}
\label{fig:convergencia}
\end{figure}

El balance numérico del residuo axial $|R_{\text{tot}} + P + qL| / |P + qL| \le 10^{-8}$ se satisface estrictamente en todas las mallas con un residuo de orden $10^{-11}\,\text{N}$[cite: 1, 30].

\section{Interpretación y Conclusiones}
En las mallas M2, M4 y M8, el desplazamiento nodal $u(L) = \SI{0.3000}{\mm}$ coincide exactamente con la solución analítica debido a la naturaleza nodal exacta de los elementos 1D bajo carga uniforme[cite: 1, 11, 22]. Sin embargo, el error entre nodos no es nulo; este disminuye cuadráticamente ($\mathcal{O}(h^2)$) al refinar la malla, como demuestra la pendiente en la Figura~\ref{fig:convergencia}[cite: 1, 11, 30]. Al imponer $u(L)=0$, el sistema se vuelve hiperestático, redistribuyendo las fuerzas a $R_{\text{izq}}=\SI{-10000}{\N}$ y $R_{\text{der}}=\SI{-30000}{\N}$[cite: 1, 12, 30]. El equilibrio estático cierra numéricamente porque el resolutor resuelve $\mathbf{K} \mathbf{u} = \mathbf{f} + \mathbf{R}$, pero esto corresponde a un problema físico completamente modificado, demostrando que la verificación de equilibrio no garantiza la validez cinemática del modelo[cite: 1, 8, 12].

\section*{Declaración de uso de IA}
Se utilizó el asistente de Inteligencia Artificial Gemini para la resolución de errores en el código de Python, la estructuración de la plantilla en LaTeX y la predicción teórica de los resultados.


%Estilo de bibliografía
\bibliographystyle{plain}

% Muestra todas las referencias del archivo .bib aunque no estén citadas explícitamente en el texto
\nocite{*}

% Nombre del archivo .bib (referencias.bib)
\bibliography{referencias}

\end{document}
